\subsection{Informed-Listener Identification Study (Pre-Registration)}
\label{sec:listener_study_prereg}

\paragraph{Status.} This subsection is the pre-registration of the informed-listener study reported in \S\ref{sec:listener_study_results} below. It is written before data collection begins (15 May 2026) and states the hypotheses, study design, sample-size target, analysis plan, and exclusion criteria so that the post-data results can be interpreted as a confirmatory rather than exploratory analysis. The companion-repository survey code, stimuli, and analysis script are at \texttt{/Volumes/MAC\_M3\_Store/Projects/listener\_study\_2026/} and were finalised prior to the first response being collected. After the data-collection period closes, the analysis script (\texttt{analysis/analyze.py}, also finalised pre-data) is executed without further modification.

\paragraph{Motivation.} Three thesis claims so far rest on metric-side evidence only: (i) the corpus mean $\mathcal{J} = 1.33$ falls below the 292c A/B within-pair family floor of $1.61$--$1.96$ (\S\ref{sec:perceptual_empirical}, \S\ref{sec:va_accuracy}), (ii) the captured corpus exhibits system-level acoustic clustering with Soundfreak tightest ($+0.79$ silhouette in hand-crafted features), Serge moderate ($+0.26$), and Buchla 200 heterogeneous-by-design ($-0.04$) per the design-philosophy interview record (\S\ref{sec:system_signatures}, Test~4), and (iii) the methodology pivot frames model evaluation as peer-distance against captured units rather than fit-quality against an abstract canonical reference (\S\ref{sec:perceptual}). The listener-side evidence at present rests on the Superbooth expert-elicitation interviews with Verbos and Buchla USA staff (\cite{verbos2026superbooth,buchla2026superbooth}). This study adds an informed-listener identification track --- methodologically distinct from a formal MUSHRA panel (which the four-layer panel-impossibility argument of \S\ref{sec:perceptual} rules out for this corpus) --- as a fourth measurement chain alongside the three metric-side chains already reported (sample-level ESR, hand-crafted Mahalanobis, CLAP).

\paragraph{Hypotheses.}
\begin{itemize}\itemsep0pt
    \item \textbf{H1 (System-character identification).} Informed listeners (modular-synth practitioners recruited via Reddit \texttt{r/modular}, Lines, Discord, and personal networks) will identify the source system of VA-rendered musical demos at above-chance accuracy. Chance baseline is $25\%$ across the four-option choice (Buchla~200 / Serge / EMS-VCS3 / Soundfreak); a result above $25\%$ supports H1.
    \item \textbf{H2 (System-clustering ordering replicates at listener level).} The per-system identification accuracy will order as Soundfreak $>$ Serge $>$ Buchla~200, replicating the silhouette-clustering pattern of Test~4 (\S\ref{sec:system_signatures}). EMS/VCS3 accuracy is exploratory (only one stimulus in the corpus). H2 is supported if the partial order $\text{acc}_\text{Soundfreak} > \text{acc}_\text{Buchla~200}$ holds (the most conservative within-pair test), and more strongly supported if the full ordering holds.
\end{itemize}
A capture-vs-VA peer-distance hypothesis was considered and intentionally excluded from this study: the capture-vs-VA distance is already quantified by the J=1.33 vs family-floor metric chain (\S\ref{sec:perceptual_empirical}); a redundant listener-side A/B of calibration sweeps would have added listening-fatigue without providing measurement-chain-independent information.

\paragraph{Study design.} Web-deployed online survey (HTML/JavaScript + Google Apps Script endpoint + Google Sheets storage). 21 trials per participant, presented in randomised order with the additional constraint that repeat-consistency stimuli must be separated by $\geq 3$ trials. Trial composition:
\begin{itemize}\itemsep0pt
    \item \textbf{17 brand-identification trials.} VA musical demos (7--12\,s each, $-23$\,dB LUFS-equivalent RMS-normalised): 9 Buchla~200 (292 LPG, 292 vintage, 281, 285, 296, 259, 265, 291 filter, 277 delay), 6 Serge (PCO, VCFQ, DUSG, wave multiplier, NTO, frequency shifter), 2 Soundfreak (filter, envelope). For each, the listener selects from \{Buchla~200 / Serge / EMS-VCS3 / Soundfreak / I don't recognise / Sounds modern-digital\} plus a 1--5 confidence rating. EMS/VCS3 remains as a foil-option in the choice menu — no source rendered cleanly enough to include as a stimulus, so any ``EMS'' response is a guaranteed brand-confusion data point.
    \item \textbf{2 modern-digital placebos.} Synthesised FM-bell arpeggio and supersaw-lead, presented in brand-ID format. Expected response is ``modern/digital'' or ``I don't recognise''; any vintage-modular response counts as an attention failure.
    \item \textbf{2 repeat-consistency trials.} Two of the brand-ID stimuli are presented a second time at a random later position. Same answer on both presentations counts as consistent.
\end{itemize}

\paragraph{Sample target and study framing.} Realistic recruitment yield from the available online modular-synth communities is $N = 10$--$15$ participants over the data-collection window. \emph{This is a pilot-scale exploratory listener study, not a confirmatory inferential one.} The sample size is too small for adequate statistical power on per-system inferential testing, and we report results accordingly: descriptive statistics and effect-size estimates with bootstrap CIs, plus non-parametric binomial-test comparisons against chance where appropriate. The study's purpose is to (a) generate listener-side evidence that complements the metric-side measurement chains, and (b) test whether the per-system ordering predicted by Test~4 (\S\ref{sec:system_signatures}) is qualitatively observed in listener data --- not to provide statistically-powered confirmatory results. A larger follow-up study is identified as future work. The convenience-sampling limitation is also acknowledged: participants are self-selected from online modular-synth communities, not a random sample of the population.

\paragraph{Analyses (pre-specified, descriptive-first given small $N$).} All analyses are descriptive in primary framing, with non-parametric significance tests where appropriate. We do NOT pre-register $p$-value thresholds for confirmatory claims; observed effects are reported with bootstrap CIs and interpreted in context.
\begin{enumerate}\itemsep0pt
    \item \textbf{Per-system identification accuracy} (H1, H2). Proportion of brand-identification trials per ground-truth system where the listener choice matched ground truth, with 95\% bootstrap CIs (200 resamples). \emph{Significance test:} binomial test against the $25\%$ chance baseline (one-tailed, $H_0$: accuracy = chance, $H_1$: accuracy > chance), applied separately per system. With $N=10$--$15$ and 7 Buchla / 4 Serge / 2 Soundfreak / 1 EMS stimuli per participant, total per-system trial counts are approximately 105/60/30/15 — the Buchla and Serge estimates are reasonably powered for the binomial test, Soundfreak and EMS are reported descriptively only.
    \item \textbf{Confusion matrix} of true-system $\times$ listener-choice. Raw counts and row-normalised proportions. Primary descriptive output of the study.
    \item \textbf{Inter-rater agreement} (descriptive). Krippendorff's $\alpha$ on categorical brand-identification responses. Reported with the standard thresholds (``acceptable'' $\geq 0.667$, ``strong'' $\geq 0.8$, ``weak'' below $0.667$). With $N=10$--$15$ this is a descriptive statistic, not powered for confirmatory testing.
    \item \textbf{Stratified accuracy} (descriptive only). By experience level, familiarity, age band, and listening device. With $N=10$--$15$ stratified cells are too small for inferential testing; reported as qualitative descriptive patterns.
    \item \textbf{Confidence-stratified accuracy} (descriptive). Per-confidence-bin (1--5) accuracy --- expected to monotonically increase if listeners are well-calibrated.
\end{enumerate}

\paragraph{Attention-check exclusion criteria (pre-specified).}
\begin{itemize}\itemsep0pt
    \item \textbf{Placebo pass criterion.} A participant ``passes'' the modern-digital placebos if at least $1$ of the $2$ placebo trials receives a ``modern/digital'' or ``I don't recognise'' response (i.e., not a vintage-modular brand claim). Participants failing this criterion are reported in a sensitivity analysis (results with and without their data) but are not the primary analysis.
    \item \textbf{Repeat-consistency criterion.} A participant ``passes'' the repeat-consistency check if at least $1$ of the $2$ repeat-pairs receives the same brand answer on both presentations. Sensitivity analysis is again reported with and without failers.
    \item \textbf{Prior-attempt criterion.} Participants who self-report having completed the study previously are excluded from the primary analysis (sensitivity analysis includes them).
    \item Per-participant trial-completion criterion: only participants who submitted all 21 trials are included.
\end{itemize}

\paragraph{Pre-stated interpretation rules (descriptive given small $N$).}
\begin{itemize}\itemsep0pt
    \item \textbf{H1 (above-chance identification).} Supported if at least one system's per-system accuracy is statistically above $25\%$ chance baseline by binomial test (one-tailed, $p < 0.05$). The systems with most stimuli (Buchla 7, Serge 4) are best-powered for this test.
    \item \textbf{H2 (ordering replicates Test~4).} Supported descriptively if the per-system accuracy ordering matches the silhouette ordering (Soundfreak $>$ Serge $>$ Buchla~200) in the attention-check-passing subset; \emph{not} formally significance-tested given the small $N$ and the descriptive-pilot framing.
    \item \textbf{If H1 and H2 are not supported:} the result is reported honestly. Two failure modes have prepared interpretations: (a) listeners cannot identify systems above chance at this stimulus set (would indicate VAs are too generic OR the demos are not representative); (b) ordering does not match Test~4 silhouettes (would indicate listener-level character does not align with the feature-space-level character --- itself an interesting finding).
\end{itemize}

\paragraph{What this study does NOT test.} The study is not a formal MUSHRA panel (the four-layer impossibility argument of \S\ref{sec:perceptual} explains why MUSHRA is inappropriate for this corpus), is not a controlled-laboratory listener experiment (online deployment with self-reported listening devices), and does not compare the deployed DSP+DE VA pipeline against alternative analog-modelling approaches (no Wright-LSTM or differentiable-solver-hybrid comparison stimuli are included --- those comparisons are reported metric-side in \S\ref{sec:cross_corpus_empirical} and \S\ref{sec:hybrid_vs_pure}). The study tests three specific listener-side complements to the existing metric-side evidence; it does not displace any of the existing thesis evidence. With $N=10$--$15$, this study is methodologically a \emph{descriptive pilot}, not a confirmatory inferential listener panel. The thesis's load-bearing empirical claims rest on the metric-side evidence (\S\ref{sec:perceptual_empirical}, \S\ref{sec:va_accuracy}, \S\ref{sec:heldout_generalisation}, \S\ref{sec:cross_corpus_empirical}); this listener pilot is a supplementary evidence type that adds qualitative human-side support, not a primary anchor.

\paragraph{Ethics and GDPR.} The study is anonymous; no personally-identifying information is collected (no name, email, IP address, browser fingerprint, cookies, exact age, location coordinates). Demographics collected are categorical bands only: age band, experience level, years-of-musical-practice band, familiarity level, listening device type, country (optional). Lawful basis is consent (Article 6(1)(a) GDPR); participants must tick an explicit consent checkbox on the welcome screen before starting. Data are stored in a Google Sheet in the researcher's Google Drive (Google's GDPR-compliant infrastructure) and will be deleted within one year of thesis submission. AAU SMC ethics-committee guidance was solicited prior to launch and the anonymous-low-risk-online-survey design is consistent with their fast-tracked-exempt criteria.

\paragraph{Reporting plan.} Results will be reported in \S\ref{sec:listener_study_results} following the structure: (a) sample description (N, demographics, attention-check pass rates), (b) primary analyses, (c) secondary analyses, (d) honest disclosure of limitations (convenience sample, uncontrolled acoustic environment, stimulus distribution imbalance across systems), (e) interpretation relative to the methodology pivot's predictions. Both the raw per-trial responses (anonymised) and the aggregated analysis outputs will be made available in the companion repository.
